\section{Discussion, Limitations, and Conclusion}
\label{sec:discussion-conclusion}
This work introduces a new lens under which circuits for modular addition may be compared.
We argue that networks trained on modular addition tend to learn a common logit manifold, and that representations at intermediate layers are dictated by the alignment of learned phases for the two inputs in the first layer preactivations.
Studying the distribution over phases, we identify that networks learn torus or vector-addition disc manifolds (known colloquially as \emph{pizzas}) in the first layer, and proceed by iteratively applying rotations and linear projections to these manifolds before ultimately arriving at a logit annulus. We argue our work recovers the universality hypothesis by showing that the counter-example architectures of \cite{zhong2023the}, and even our torus-learning MLP-Concat architecture all use the same torus to logits map. The transformers and MLP-Add architecture simply learn lower rank manifolds of the same class of torus manifolds. This follows rigorously and directly from our closed-form equations for the torus and pizza discs: the pizza discs are a linear projection of the torus. Thus, we posit ``is the nature of universality that DNNs recover either a universal manifold, or linear projections of it in order to fit data?'' 

Notably, this conclusion was reached by examining networks from a different perspective, which allowed us to mathematically predict the structure of representations, and also to test these predictions precisely at large scale using tools from topological data analysis. In particular, our work identifies that the joint distribution over simple neuron phases both determines the geometry of representations, and permits efficient quantitative evaluation.

A notable limitation of our work is that, in its present form, it treats only circuits in modular addition.
To interpret large-scale neural networks, it will be necessary to derive strategies for characterizing circuits across domains.
While our precise methodology is particular to the analysis of modular addition networks, we believe the underlying strategy of isolating the degrees of freedom in the geometry of hidden representations and estimating the statistics of their corresponding learned values can be a generalizable approach to characterize circuits in more complex domains.

Beyond this, our results also suggest that the universality property is likely to hold in modular addition.
It also draws connections between universality and the manifold hypothesis, given the correspondence we see between the low-dimensional representations formed throughout the layers of the networks we observed.
Ultimately, we believe the tools we introduced in this work can inform formal hypotheses about universality and manifold continuity that can be assessed quantitatively.